\documentclass[aspectratio=169,10pt,spanish]{beamer}

\input{../../../_preambulo_beamer.tex}
\input{../../../_comandos_beamer.tex}

\selectlanguage{spanish}

\title{MA1001B - Modelación Estadística para la Toma de Decisiones}
\subtitle{Lección 08: Reglas de probabilidad y tablas de contingencia}
\author{}
\institute{Tecnol\'ogico de Monterrey}
\date{}

\begin{document}

\begin{frame}[plain]
  \vspace{1.2cm}
  {\Large\bfseries MA1001B - Modelación Estadística para la Toma de Decisiones\par}
  \vspace{0.7cm}
  {\large Lección 08: Reglas de probabilidad y tablas de contingencia\par}
  \vspace{0.7cm}
  {\color{TecRojo}\rule{\textwidth}{1.2pt}}
  \vspace{0.8cm}

  Facultad MA1001B\\[0.3cm]
  Periodo\\[0.3cm]
  Tecnol\'ogico de Monterrey
\end{frame}

\begin{frame}{La pregunta de probabilidad}
  \begin{alertblock}{Pregunta guía}
    ¿Cómo puede una tabla de contingencia sostener probabilidades correctas de
    Y, O y NO sin contar dos veces ni asumir independencia?
  \end{alertblock}

  \vspace{0.4em}
  Al final de esta lección, usted puede:
  \begin{itemize}
    \item leer probabilidades marginales y conjuntas en una tabla de doble entrada;
    \item aplicar las reglas de complemento y de adición general;
    \item identificar el traslape en una probabilidad O;
    \item explicar por qué multiplicar probabilidades marginales exige evidencia
      de independencia.
  \end{itemize}

  \vfill
  \scriptsize Todos los registros de pedidos y procesos operativos son sintéticos.
\end{frame}

\begin{frame}{Una tabla de contingencia organiza resultados conjuntos}
  \small
  Sea \(M\) revisión manual y \(L\) entrega retrasada.

  \begin{center}
    \begin{tabular}{lrrr}
      \toprule
      \textbf{Proceso de enrutamiento} & \textbf{Retrasado} & \textbf{A tiempo} & \textbf{Total} \\
      \midrule
      Revisión manual & 48 & 72 & 120 \\
      Automatizado & 32 & 248 & 280 \\
      \midrule
      Total & 80 & 320 & 400 \\
      \bottomrule
    \end{tabular}
  \end{center}

  \vspace{0.6em}
  \begin{columns}[T]
    \begin{column}{0.48\textwidth}
      \textbf{Probabilidades marginales}
      \[
        P(M)=\frac{120}{400}=0.30,\quad
        P(L)=\frac{80}{400}=0.20.
      \]
    \end{column}
    \begin{column}{0.48\textwidth}
      \textbf{Probabilidad conjunta}
      \[
        P(M\cap L)=\frac{48}{400}=0.12.
      \]
    \end{column}
  \end{columns}
\end{frame}

\begin{frame}{Paso 1: Leer celdas y márgenes antes de aplicar reglas}
  \begin{columns}[T]
    \begin{column}{0.36\textwidth}
      \small
      \begin{block}{Salida verificada}
        Pedidos sintéticos: 400\\
        \(P(M)=0.30\)\\
        \(P(L)=0.20\)\\
        \(P(M\cap L)=0.12\)
      \end{block}

      Una celda representa dos condiciones a la vez. Los totales de fila y de
      columna representan una condición sin importar la otra variable.
    \end{column}
    \begin{column}{0.64\textwidth}
      \centering
      \includegraphics[width=0.88\textwidth]{../figuras/probabilidad_01_tabla.png}
    \end{column}
  \end{columns}
\end{frame}

\begin{frame}{Reglas de complemento y adición}
  \small
  La regla del complemento da cuenta de todo lo que queda fuera de un evento:
  \[
    P(M^c)=1-P(M)=1-0.30=0.70.
  \]

  La regla general de adición elimina los resultados compartidos contados dos
  veces:
  \[
    P(M\cup L)=P(M)+P(L)-P(M\cap L).
  \]
  \[
    P(M\cup L)=0.30+0.20-0.12=0.38.
  \]

  \begin{alertblock}{Trampa frecuente}
    La suma ingenua 0.50 incluye dos veces los 48 pedidos de revisión y
    retraso. La suma simple aplica solo cuando los eventos son mutuamente
    excluyentes.
  \end{alertblock}
\end{frame}

\begin{frame}{Paso 2: La regla de adición coincide con el conteo directo}
  \begin{columns}[T]
    \begin{column}{0.35\textwidth}
      \small
      \begin{block}{Comparación verificada}
        Suma ingenua: 0.50\\
        Traslape: 0.12\\
        Unión corregida: 0.38\\
        Conteo directo: 0.38
      \end{block}

      Contar de forma directa todos los pedidos en \(M\cup L\) da
      \[
        \frac{152}{400}=0.38.
      \]
    \end{column}
    \begin{column}{0.65\textwidth}
      \centering
      \includegraphics[width=\textwidth]{../figuras/probabilidad_02_adicion.png}
    \end{column}
  \end{columns}
\end{frame}

\begin{frame}{Paso 3: El atajo de multiplicación tiene una condición}
  \begin{columns}[T]
    \begin{column}{0.38\textwidth}
      \small
      Regla general:
      \[
        P(M\cap L)=P(M)P(L\mid M).
      \]
      La tabla da \(P(L\mid M)=48/120=0.40\), de modo que
      \[
        0.30(0.40)=0.12.
      \]

      Multiplicar las marginales da \(0.30(0.20)=0.06\), que contradice la
      intersección observada.

      \textbf{Límite:} la independencia no se ha justificado.
    \end{column}
    \begin{column}{0.62\textwidth}
      \centering
      \includegraphics[width=\textwidth]{../figuras/probabilidad_03_limite.png}
    \end{column}
  \end{columns}
\end{frame}

\begin{frame}{Verificación de conceptos}
  Use la tabla de 400 pedidos de esta lección.
  \begin{enumerate}
    \item Encuentre \(P(\text{A tiempo})\).
    \item Encuentre \(P(\text{Automatizado O Retrasado})\) con la regla de
      adición.
    \item ¿Por qué \(P(M)+P(L)=0.50\) no es la probabilidad de \(M\cup L\)?
    \item ¿Qué evidencia muestra que \(P(M)P(L)\) no es válido para esta
      intersección?
  \end{enumerate}

  \vspace{0.4em}
  \begin{block}{Conexión con la Lección 09}
    La siguiente lección desarrolla probabilidad condicional e independencia:
    cómo un espacio muestral reducido cambia la probabilidad y cómo se prueba
    la independencia.
  \end{block}

  \vfill
  \scriptsize Fuente: OpenStax, \textit{Introductory Business Statistics 2e},
  Capítulo 3, Secciones 3.1 y 3.3. CC BY-NC-SA 4.0.
\end{frame}

\appendix



\end{document}
